\documentclass[11pt]{article}
\usepackage{amsmath,amssymb}
\usepackage[margin=1in]{geometry}
\title{Exact limit of the symmetric probe inequalities}
\author{}
\date{}
\begin{document}
\maketitle

This calculation optimizes only the following arithmetic scheme.
Let $\epsilon=7/8-B\geq0$ and let $t\geq0$ change the lengths by
\[
 \ell=\frac16+t,\qquad
 \ell_x=\frac{17}{48}-\frac t2,\qquad
 \ell_y=\frac{23}{48}-\frac t2,\qquad
 h=\frac{13}{16}+\frac{3t}{2}.
\]
For $0\leq\delta\leq3/4$ and $0\leq y\leq1/2$, put
\[
 P=\frac{7+18y+8y^2}{9},\quad
 J=\left(\frac56-\delta\right)\frac{37+34y}{18}+\delta P,
 \quad L=\left(\frac56-\delta\right)\delta P,
\]
\[
 R=1-\delta+\frac{L}{2J},\quad
 A=\frac{1+(3+8y)\delta}{48}-\frac{13L}{32J},\quad
 S=-\frac12+\delta+\delta\left(\frac12-y\right)+\frac{3R}{2}.
\]
The constraints are
\[
 \epsilon<\frac t4,
 \qquad -A+\epsilon+tS<0\quad\hbox{throughout the rectangle}.
\]
They use the comparison exponent at $B$. Additional analytic constraints
and a comparison exponent depending on an individual zero are not free
parameters in this optimization.

Since $J>0$, $1-\delta\leq R\leq17/12$. Hence
$5/8\leq S\leq11/4$. Taking $t$ down to $4\epsilon$ reduces the
supremum problem to the minimum of $A/(1+4S)$.

Define $q(e)=49-1144080e-18604800e^2$. Its unique positive root is
\[
 e_* = \frac{98}{1144080+
   \sqrt{1144080^2+3646540800}}.
\]
The root is positive and is less than $1/20000$.

Here is a global certificate. Write
\[
 5184J\bigl(A-e(1+4S)\bigr)=a_e(y)\delta^2+b_e(y)\delta+c_e(y),
\]
where
\begin{align*}
 a_e(y)&=1224+12096e+(3144+4608e)y\\
 &\quad +(2256+16128e)y^2+(768+18432e)y^3,\\
 b_e(y)&=-948+23040e+(-1508+6720e)y+(-104-1920e)y^2,\\
 c_e(y)&=185-44400e+(170-40800e)y.
\end{align*}
Direct polynomial expansion gives
\begin{align*}
 4a_ec_e-b_e^2
 &=144q(e)\\
 &\quad +(299712-664266240e-3102105600e^2)y\\
 &\quad +(1336112-877278720e-3573043200e^2)y^2\\
 &\quad +(1788736-484362240e-5879808000e^2)y^3\\
 &\quad +(511424-113203200e-3011788800e^2)y^4.
\end{align*}
For $0\leq e\leq1/20000$, the last four coefficients have the
respective positive lower bounds
\[
 \frac{4163920824}{15625},\quad
 \frac{20191236428}{15625},\quad
 \frac{5514072464}{3125},\quad
 \frac{7902442352}{15625}.
\]
At $e=e_*$ the constant coefficient vanishes and all other coefficients
are positive. Since $a_e(y)>0$ for $e,y\geq0$, the identity
\[
 4a_e\bigl(a_e\delta^2+b_e\delta+c_e\bigr)
 =(2a_e\delta+b_e)^2+(4a_ec_e-b_e^2)
\]
proves $A\geq e_*(1+4S)$ on the whole rectangle. Equality holds only at
\[
 y=0,\qquad
 \delta=\delta_*:=\frac{948-23040e_*}{2448+24192e_*}.
\]
This point belongs to the prescribed rectangle.

For $0\leq e<e_*$, choose $t=4e+(e_*-e)/4$. The low inequality is
strict. The high margin satisfies
\begin{align*}
 A-e-tS
 &\geq (e_*-e)\left(1+\frac{15S}{4}\right)>0.
\end{align*}
For $e\geq e_*$, the strict low inequality forces $t>4e$. Evaluating
at $(\delta_*,0)$ gives
\[
 -A+e+tS=(e-e_*)(1+4S)+(t-4e)S>0.
\]
Thus the strict arithmetic scheme is feasible exactly when $e<e_*$.
The root is its supremum and is not itself feasible.

The reproducible calculation is
\texttt{uv run research/method\_limit\_symmetric.py}.
Its exact rational checkpoint is
\[
 \frac{21399692313119}{500000000000000000}
 <e_*<
 \frac{42799384626239}{1000000000000000000}.
\]
The endpoint polynomial has the exact signs
\[
 q\left(\frac{21399692313119}{500000000000000000}\right)
 =\frac{35831290538249775973}{39062500000000000000000000000000}>0,
\]
\[
 q\left(\frac{42799384626239}{1000000000000000000}\right)
 =-\frac{35686173469217846747}{156250000000000000000000000000000}<0.
\]
Numerically,
\[
 e_* =0.0000427993846262388006485268367\ldots,
 \qquad \delta_* =0.3866885309739388256550471335\ldots,
\]
\[
 B_* =\frac78-e_*=0.8749572006153737611993514732\ldots.
\]

This is a limit of the stated inequalities, not of every probe method.
Although the strict margins vanish at the limit, a proved family of
analytic bounds for all $e<e_*$ could exclude every zero with real part
greater than $B_*$. That closure argument is separate from this
arithmetic certificate. Changing the comparison exponent to depend on
the zero, or adding another independent geometric parameter, requires
stating and optimizing those altered inequalities.

\end{document}
